\section{Solution}

We propose a progressive fine-tuning pipeline that addresses the
compositional reasoning gap exposed by our baseline evaluation.
The pipeline adapts a frozen VLM through three cumulative stages, each
targeting a specific failure mode: format compliance, structured reasoning,
and distractor discrimination.

\subsection{Base Model and Adaptation}

We adopt InternVL2.5-8B~\cite{chen2024internvl25} as the base VLM and
adapt it via LoRA~\cite{hu2022lora} applied to both the language model and
the vision--language projector, yielding 19M trainable parameters out of
${\sim}8$B total (0.23\%).
The vision encoder remains frozen.
Each video is represented as 8 uniformly-sampled frames producing
$2{,}048$ visual tokens within a $3{,}072$-token sequence budget.
Full hyperparameters are listed in the appendix.

\subsection{Phase~1: Supervised Fine-Tuning}
\label{sec:sft}

The base model's zero-shot responses are free-form and frequently fail to
commit to a single MCQ option.
SFT teaches the model to map video frames and a multiple-choice prompt to a
lettered answer, using standard causal language modeling loss masked to the
response tokens.
We allocate 20\% of videos to training and 80\% to test, with a
group-aware split that keeps clips from the same parent video in the same
partition to prevent data leakage.
Annotation text is withheld from the prompt so the model must ground its
answer in visual evidence.

\subsection{Phase~2--3: Chain-of-Thought Distillation and Training}
\label{sec:cot}

Our baselines show that accuracy drops monotonically with compositional
demand: models that score 50\%+ on single-attribute questions fall below
30\% when asked to jointly identify aggressor, action, and victim.
Direct answer supervision does not teach the intermediate reasoning steps
needed to compose these attributes.

To address this, we distill chain-of-thought
reasoning~\cite{wei2022chain} from the base model via a two-stage
generation process: the teacher first produces a vision-grounded reasoning
chain, then refines it against the ground-truth annotation.
A quality filter ensures every chain commits to the correct MCQ letter.
The resulting chains are used as supervision for compound and complex
questions, while simple questions retain direct answers.
Phase~3 continues training the Phase~1 adapter on this mixed dataset,
adding structured reasoning without losing format compliance.

\subsection{Phase~4--5: Taxonomy-Guided Preference Optimization}
\label{sec:adpo}

SFT and CoT teach the model \emph{what} to answer but not \emph{what to
avoid}.
To sharpen discrimination against plausible distractors, we apply anchored
direct preference optimization
(ADPO)~\cite{zhang2025adpo,rafailov2023direct}.

\paragraph{Taxonomy-graded pair construction.}
Each distractor in our benchmark carries a hardness label from a 10-level
taxonomy ranging from \emph{role reversal} (hardest --- aggressor and victim
swapped, both visible) to \emph{cross-video} (easiest --- neither person
appears).
For each training question, we pair the correct answer against the single
hardest-ranked distractor, concentrating the preference signal on the
distinctions models struggle with most.

\paragraph{Anchored DPO loss.}
Let $\pi_\theta$ denote the policy and $\pi_{\mathrm{ref}}$ the reference
(same weights with LoRA disabled).
For chosen response $y^+$ and rejected $y^-$ conditioned on prompt $x$:
%
\begin{equation}
  \Delta^{\pm} = \log \pi_\theta(y^{\pm} \mid x) - \log \pi_{\mathrm{ref}}(y^{\pm} \mid x)
\end{equation}
%
\begin{equation}
  \mathcal{L} = \underbrace{-\log\sigma\!\bigl(\beta(\Delta^+ - \Delta^-)\bigr)}_{\text{preference}}
              + \alpha\underbrace{\bigl((\Delta^+)^2 + (\Delta^-)^2\bigr)}_{\text{anchor}}
\end{equation}
%
The anchor term ($\alpha = 0.5$) prevents unbounded policy drift while the
preference term ($\beta = 0.1$) steers the model toward the correct answer.
Setting $\alpha = 0$ recovers vanilla DPO.
